\documentclass[a4paper,12pt,twoside,openany]{book}
% Paquete para determinar la codificación del archivo. Las opciones usuales son utf8 y latin1
\usepackage[utf8]{inputenc}
% Paquete para personalizar el idioma.
\usepackage[english,spanish]{babel}
% Paquetes de fuentes
\usepackage{mathptmx}
\usepackage{helvet}
% Paquete para hiperenlaces
\usepackage[colorlinks]{hyperref}
\hypersetup{linkcolor=Blue,citecolor=Green}
% Paquete que contiene las definiciones de la plantilla
\usepackage{tfg-estilo}

% Paquete para acrónimos
\usepackage[acronym]{glossaries}
%% Definición de acrónimos
%% Uso: \gls{TFG} para la primera mención (expandido)
%%      \glspl{TFG} para plural
%%      \glsxtrshort{TFG} solo la abreviatura

\newacronym{ML}{ML}{Machine Learning}
\newacronym{SL}{SL}{Supervised Learning}
\newacronym{UL}{UL}{Unsupervised Learning}
\newacronym{RL}{RL}{Reinforcement Learning}
\newacronym{SSL}{SSL}{Semi-Supervised Learning}
\newacronym{SUL}{SUL}{Semi-Unsupervised Learning}
\newacronym{ERM}{ERM}{Empirical Risk Minimization}
\newacronym{RRM}{RRM}{Regularized Risk Minimization}
\newacronym{ST}{ST}{Self-Training}
\newacronym{CT}{CT}{Co-training}
\newacronym{CoReg}{CoReg}{Co-regularization}
\newacronym{PCA}{PCA}{Principal Component Analysis}
\newacronym{AE}{AE}{Autoencoder}
\newacronym{EM}{EM}{Expectation-Maximization}
\newacronym{SVM}{SVM}{Support Vector Machine}
\newacronym{TSVM}{TSVM}{Transductive SVM}
\newacronym{S3VM}{S3VM}{Semi-Supervised Support Vector Machine}
\newacronym{NN}{NN}{Neural Network}
\newacronym{DNN}{DNN}{Deep Neural Network}
\newacronym{FNN}{FNN}{Feedforward Neural Network}

%% Añade aquí más acrónimos en el formato:
%% \newacronym{acrónimo}{forma corta}{forma expandida}

\makenoidxglossaries

{\theorembodyfont{\itshape}
\newtheorem{theo}{Teorema}[chapter]
\newtheorem{prop}[theo]{Proposición}
\newtheorem{coro}[theo]{Corolario}
\newtheorem{lemm}[theo]{Lema}
} {\theorembodyfont{\upshape}
\newtheorem{rema}[theo]{Observación}
\newtheorem{defi}[theo]{Definición}
\newtheorem{algo}[theo]{Algoritmo}
\newtheorem{example}[theo]{Ejemplo}
}
\def\theequation{\thesection.\arabic{equation}}
\def\thetheo{\thesection.\arabic{theo}}
\newenvironment{proof}{\ignorespaces\par\noindent\textbf{Demostración.}\quad}{\hfill$\Box$\par}

\begin{document}

\input{portada.tex}

\input{agradecimientos.tex}

\frontmatter
\pagestyle{fancyplain}

\tableofcontents

\chapter*{Abstract}
\makeatletter\@mkboth{Abstract}{Abstract}\makeatother
\addcontentsline{toc}{chapter}{Abstract}


In the last decade, machine learning has emerged as the driving force behind revolutionary advances in areas such as computer vision, natural language processing, and autonomous robotics. However, the success of the most powerful models, particularly in the supervised learning paradigm, critically depends on the availability of vast labeled datasets. The labeling process, often performed by human experts, constitutes a logistical and economic bottleneck, being in many cases a laborious, slow, and error-prone task. In this scenario, \acrfull{SSL} emerges as an essential discipline that seeks to maximize classifier performance using a small set of labeled data in combination with a large amount of unlabeled data.

This Bachelor's thesis addresses the problem of semi-supervised classification from an analytical perspective, analyzing how knowledge of the geometry of the input space allows refining decision boundaries. Unlike other more empirical approaches, this study focuses on the mathematical foundations that justify the use of unlabeled data. For these data to add value, certain regularity assumptions must be met. In this regard, the three fundamental assumptions of \acrshort{SSL} are formalized and discussed: the smoothness assumption, which postulates that nearby points in high-density regions should share the same label; the cluster or low-density assumption, which suggests that the decision boundary should pass through regions where data are scarce; and the manifold assumption, which assumes that high-dimensional data concentrate on a lower-dimensional structure. Understanding these assumptions is crucial, as if the data do not satisfy them, including unlabeled examples can introduce noise and degrade the performance of the original model.

The central contribution of this work is the proposal of a structured taxonomy to organize the vast literature of \acrshort{SSL}, inspired by recent surveys, especially \citet{van2020survey}, and adapted with motivated modifications to the aims of this study. Through a didactic yet rigorous approach, algorithms are classified according to their mathematical nature and how they incorporate unsupervised information. The work not only compiles scattered information but synthesizes it through a comparative analysis of the loss functions and regularization terms specific to each paradigm. The goal is to provide the reader with a conceptual map that allows discerning between inductive methods, designed to learn a generalizable decision rule for new data, and transductive methods, whose purpose is to infer labels exclusively for a previously defined set of unlabeled points.

The proposed taxonomy is developed in detail through four major blocks. First, wrapper methods, such as self-training and disagreement-based methods, are analyzed, which use supervised models iteratively to generate pseudo-labels. Second, unsupervised preprocessing is studied, where feature extraction and clustering techniques are employed to transform the representation space before applying the classifier. The third block, of particular mathematical interest, analyzes intrinsically semi-supervised methods. Here, generative models that use the EM algorithm to optimize mixtures of distributions and discriminative methods such as Semi-supervised Support Vector Machines (S3VM), which seek to maximize the margin in the presence of unlabeled data, are explored in depth. Finally, graph-based methods are explored, analyzing the construction of adjacency matrices and the use of harmonic functions and mincut for label propagation.

In conclusion, this work establishes a solid theoretical framework that clarifies the mathematical foundations of current \acrshort{SSL} techniques. The systematic review of algorithms and their properties highlights that the success of \acrshort{SSL} does not depend on the amount of data per se, but on the correct alignment between model assumptions and the underlying geometric structure of the data. This study lays the groundwork for future lines of research, such as the empirical evaluation of these taxonomies in large language models or advanced vision systems, ensuring a smooth transition between mathematical theory and practical implementation.

\chapter*{Resumen}
\makeatletter\@mkboth{Resumen}{Resumen}\makeatother
\addcontentsline{toc}{chapter}{Resumen}

En la última década, el aprendizaje computacional se ha consolidado como el motor fundamental detrás de avances revolucionarios en áreas como la visión artificial, el procesamiento del lenguaje natural y la robótica autónoma. No obstante, el éxito de los modelos más potentes, particularmente en el paradigma del aprendizaje supervisado, depende críticamente de la disponibilidad de vastos conjuntos de datos etiquetados. El proceso de etiquetado, a menudo realizado por expertos humanos, constituye un cuello de botella logístico y económico, siendo en muchos casos una tarea laboriosa, lenta y propensa a errores. Ante este escenario, el aprendizaje semi-supervisado (SSL, por sus siglas en inglés) se presenta como una disciplina esencial que busca maximizar el rendimiento de los clasificadores utilizando un conjunto reducido de datos etiquetados en combinación con una gran cantidad de datos no etiquetados.

Este trabajo aborda el problema de la clasificación semi-supervisada desde una perspectiva analítica, analizando cómo el conocimiento de la geometría del espacio de entrada permite refinar las fronteras de decisión. A diferencia de otros enfoques más empíricos, este estudio se centra en las bases matemáticas que justifican el uso de datos sin etiquetas. Para que estos datos aporten valor, deben cumplirse ciertas hipótesis de regularidad. En este sentido, se formalizan y discuten las tres suposiciones fundamentales del SSL: la suposición de suavidad, que postula que puntos cercanos en regiones de alta densidad deben compartir la misma etiqueta; la suposición de clúster o de baja densidad, que sugiere que la frontera de decisión debe pasar por regiones donde los datos son escasos; y la suposición de variedad, que asume que los datos de alta dimensión se concentran en una estructura de menor dimensión. Comprender estas hipótesis es crucial, ya que si los datos no las satisfacen, la inclusión de ejemplos no etiquetados puede introducir ruido y degradar el rendimiento del modelo original.

La contribución central de este trabajo es la propuesta de una taxonomía estructurada que organiza la vasta literatura del aprendizaje semi-supervisado, inspirada en revisiones recientes (en particular, van Engelen y Hoos, 2020) y adaptada con modificaciones motivadas a los objetivos del estudio. A través de un enfoque didáctico pero riguroso, se clasifican los algoritmos atendiendo a su naturaleza matemática y a la forma en que incorporan la información no supervisada. El trabajo no solo recopila información dispersa, sino que la sintetiza mediante un análisis comparativo de las funciones de pérdida y los términos de regularización específicos de cada paradigma. El objetivo es proporcionar al lector un mapa conceptual que permita discernir entre métodos inductivos, diseñados para aprender una regla de decisión generalizable a nuevos datos, y métodos transductivos, cuyo propósito es inferir etiquetas exclusivamente para un conjunto de puntos no etiquetados previamente definidos.

La taxonomía propuesta se desarrolla detalladamente a través de cuatro grandes bloques. En primer lugar, se analizan los métodos envolventes, como el auto-entrenamiento y los métodos basados en desacuerdo (co-regularización y múltiples vistas), que utilizan modelos supervisados de forma iterativa para generar pseudo-etiquetas. En segundo lugar, se estudia el preprocesado no supervisado, donde se emplean técnicas de extracción de características y clustering para transformar el espacio de representación antes de aplicar el clasificador. El tercer bloque, de especial interés matemático, analiza los métodos intrínsecamente semi-supervisados. Aquí se profundiza en modelos generativos que utilizan el algoritmo EM para optimizar mezclas de distribuciones, y en métodos discriminativos como las Máquinas de Vectores de Soporte Semi-supervisadas (S3VM), que buscan maximizar el margen en presencia de datos no etiquetados. Finalmente, se exploran los métodos basados en grafos, analizando la construcción de matrices de adyacencia y el uso de funciones armónicas y cortes mínimos (mincut) para la propagación de etiquetas.

En conclusión, este trabajo establece un marco teórico sólido que clarifica los fundamentos matemáticos de las técnicas actuales de aprendizaje semi-supervisado. La revisión sistemática de los algoritmos y sus propiedades pone de manifiesto que el éxito del SSL no depende de la cantidad de datos per se, sino de la correcta alineación entre las suposiciones del modelo y la estructura geométrica subyacente de los datos. Este estudio sienta las bases para futuras líneas de investigación, como la evaluación empírica de estas taxonomías en grandes modelos de lenguaje o sistemas de visión avanzados, garantizando una transición fluida entre la teoría matemática y su implementación práctica.

\mainmatter

\input{capitulo1.tex}

\input{capitulo2.tex}

\input{capitulo3.tex}

\chapter{Metodos envolventes}
% En este capítulo vamos a proceder con el desarrollo teórico de los diferentes métodos o familias de métodos que se han presentado en la taxonomía. En cada sección se hará una breve introducción a la categoría dentro de la taxonomía, recordando sus aspectos principales que la diferencian de las demás, y posteriormente se hará una un estudio de los métodos más representativos de cada categoría, explicando su funcionamiento.
La familia de métodos envolventes es de las más viejas y ampliamente conocidas dentro del aprendizaje semi-supervisado \cite{zhou2010semi}. Como se ha explicado a la hora de presentar esta categoría en la taxonomía (Sección \ref{sec:div-uso-etiquetas}), estos métodos se basan en la idea del \textit{pseudoetiquetado}, un procedimiento iterativo donde cada iteración consiste en dos pasos: primero se entrena a uno o varios clasificadores supervisados con los datos etiquetados originales y posiblemente con datos etiquetados de iteraciones anteriores, y luego se utilizan estos clasificadores para predecir etiquetas para los datos no etiquetados, las predicciones de mayor confianza son entonces añadidas al conjunto de datos etiquetados, conviertiendose en datos pseudoetiquetados. Este proceso se repite hasta que se cumpla un criterio de parada, como por ejemplo un número máximo de iteraciones o una mejora mínima en el rendimiento del modelo.

Esta familia de métodos tiene varias peculiaridades que vamos a exponer a continuación:

\begin{itemize}
    \item En primer lugar, tienen una propiedad clave que es la de poder ser aplicados a cualquier, o casi cualquier, modelo supervisado subyacente. Esto les otorga de una gran flexibilidad pues introducen el uso de los datos no etiquetados en el proceso de entrenamiento de manera sencilla y directa. 
    \item Por otro lado, cabe recalcar que no se basan por si mismos en ningún supuesto específico sobre la distrbución de los datos o las relaciones entre los datos etiquetados y no etiquetados, si no que dependen de los supuestos que tengan los modelos supervisados subyacentes. Esto les otroga gran versatilidad, y los convierte en candidatos idoneos para mejorar el rendimiento en problemas en los que algún algoritmo supervisado ya trabaja bien. 
    \item Sin embargo tienen una suposición importante: los métodos envolventes asumen que las predicciones, al menos la de mayor confianza, son correctas. En el caso de que esto no se cumpla, se puede producir un efecto de arrastre donde los errores se van acumulando a lo largo de las iteraciones, lo que puede resultar en una degradación del rendimiento del modelo \cite{zhu2009introduction}.
\end{itemize}

Se van a presentar a continuación tres métodos o familias de métodos representativas de los métodos envolventes: el auto-entrenamiento, los métodos basados en desacuerdo y los métodos basados en boosting.

\section{Auto-entrenamiento}

El auto-entrenamiento es la versión más simple de los métodos envolventes y una de las más antiguas \cite{yarowsky1995unsupervised}. En este método se toma la idea del pseudoetiquetado y se aplica a un único clasificador supervisado, que se entrena con los datos etiquetados originales y luego se utiliza para predecir etiquetas para los datos no etiquetados, usando las predicciones de mayor confianza reforzar el entrenamiento en la siguiente iteración.

A pesar de su simplicidad, se pueden realizar varias decisiones de diseño a la hora de implementar un método de auto-entrenamiento: el criterio para seleccionar las predicciones de mayor confianza, el número de predicciones a añadir en cada iteración o el criterio de parada. Estas decisiones pueden tener un impacto significativo en el rendimiento del método y pueden ser adaptadas a las características específicas del problema y los datos \cite{van2020survey}.

Se expone a continuación un algortitmo de auto-entrenamiento donde se pueden observar estas decisiones de diseño:

\begin{algo}
    \textbf{Auto-entrenamiento}
    \\Entrada: Conjunto de datos etiquetados $L = \{(\mathbf{x}_i, y_i)\}_{i=1}^l$, conjunto de datos no etiquetados $U = \{\mathbf{x}_j\}_{j=l+1}^{l+u}$, clasificador supervisado $f$, criterio de selección de predicciones de mayor confianza $C$, criterio de parada $S$.
    \\Salida: Clasificador entrenado $f$.
    \\Procedimiento:
    \\Mientras no se cumpla el criterio de parada $S$:
    \begin{enumerate}
        \item Entrenar el clasificador supervisado $f$ con los datos etiquetados $L$.
        \item Predecir etiquetas para los datos no etiquetados $U$ utilizando el clasificador supervisado $f$, obteniendo un conjunto de predicciones $P$.
        \item Seleccionar un subconjunto de predicciones $P_{conf}$ de $P$ utilizando el criterio de selección de predicciones de mayor confianza $C$.
        \item Añadir las predicciones seleccionadas $P_{conf}$ al conjunto de datos etiquetados $L$, convirtiéndolas en datos pseudoetiquetados.
        \item Eliminar las predicciones seleccionadas $P_{conf}$ del conjunto de datos no etiquetados $U$.
    \end{enumerate}
\end{algo}

\section{basados en desacuerdo}

Los métodos basados en desacuerdo son la extensión natural del auto-entrenamiento a múltiples clasificadores supervisados. Su nombre proviene del hecho de que en cada iteración del pseudoetiquetado, una vez entrenados los clasificadores supervisados, estos no necesariamente estan de acuerdo en sus predicciones para los datos no etiquetados, y se pueden aprovechar estas diferencias para seleccionar las predicciones que se van a añadir al conjunto de datos etiquetados. Por ejemplo, se pueden seleccionar las predicciones donde los clasificadores están de acuerdo, o por el contrario, predicciones donde algun clasificador tiene mucha confianza en su predicción, dependiendo de la estrategia que se quiera seguir \cite{zhou2010semi}.

Los métodos basados en desacuerdo utilizan los datos no etiquetados como una especie de puente para intercambiar información entre los modelos a entrenar. El origen de esta discrepancia en la información y la manera en la que esta es propagada divide estos métodos en varias familias que se presentan a continuación.

\subsection{Multiples vistas}

Una de las primeras ideas que dieron forma a este tipo de métodos fue la de descomponer el dominio de las instancias $\mathcal{X}$ en lo que se conoce como vistas, es decir, subconjuntos disjuntos de las características, y entrenar un clasificador supervisado para cada vista, de esta manera cada clasificador tiene una información diferente sobre el problema y se pueden aprovechar estas diferencias para mejorar el rendimiento del modelo \cite{zhou2010semi}. 

El primero método de este tipo fue el co-training presentado en \textit{Combining labeled and unlabeled data with co-training, Blum and Mitchell, 1998} \cite{blum1998combining} donde se divide el dominio en dos vistas $\mathcal{X} = \mathcal{X}_1 \cup \mathcal{X}_2$ de tal manera que cada vector de características $\mathbf{x}$ se puede representar como $\mathbf{x} = (\mathbf{x}_1, \mathbf{x}_2)$ con $\mathbf{x}_1 \in \mathcal{X}_1$ y $\mathbf{x}_2 \in \mathcal{X}_2$. El algoritmo de co-training consiste en entrenar dos clasificadores supervisados $f_1$ y $f_2$ con las vistas $\mathcal{X}_1$ y $\mathcal{X}_2$ respectivamente, y luego utilizar cada clasificador para predecir etiquetas para los datos no etiquetados, añadiendo las predicciones de mayor confianza de cada clasificador al conjunto de datos etiquetados del otro clasificador. Este proceso se repite iterativamente hasta que se cumpla un criterio de parada.

Blum y Mitchell hacen dos suposiciones clave para que co-training funcione correctamente: la primera es que las vistas son condicionalmente independientes dado la etiqueta, es decir, $P(\mathbf{x}_1, \mathbf{x}_2 \mid y) = P(\mathbf{x}_1 \mid y) P(\mathbf{x}_2 \mid y)$, lo que implica que cada vista proporciona información complementaria sobre la etiqueta. La segunda suposición es que cada vista es suficiente para predecir la etiqueta, es decir, $P(y \mid \mathbf{x}_1) = P(y \mid \mathbf{x}_2) = P(y \mid \mathbf{x})$, lo que implica que cada vista por si sola es capaz de predecir la etiqueta con un rendimiento razonable \cite{zhu2009introduction}. 

Sin embargo, se ha desmostrado que el los métodos de multiples vistas pueden funcionar con suposiciones más débiles en la independencia \cite{abney2002bootstrapping}. Por ejemplo, la suposición de expansión \cite{balcan2004co} relaja la independencia condicional, exigiendo simplemente que los ejemplos no estén excesivamente correlacionados y que el clasificador no cometa errores con alta confianza. Desde un punto de vista de teoría de aprendizaje, esto garantiza que el error de generalización se reduce al forzar el acuerdo entre las vistas.

\subsection{Una única vista}

Sin embargo, hay una suposición aun mayor en co-training y los métodos basados en multiples vistas, y es que en muchos casos del mundo real no existe una división natural del dominio en vistas. Es aquí donde surgen los métodos de una única vista, estos buscan generar múltiples clasificadores supervisados a partir de un dominio que en su origen no tiene una división en vistas.

Una primera idea puede ser la de generar las vistas artificialmente, y usar estas vistas para entrenar los diferentes clasificadores supervisados \cite{wang2008random}. Otras formas más naturales de entrenar múltiples clasificadores supervisados con una única vista, pero introduciendo alguna fuente de aleatoriedad para que cada clasificador tenga una información diferente sobre el problema, pueden ser utilizar diferentes parametros para un mismo modelo \cite{wang2007analyzing} o utilizar directamente diferentes modelos supervisados \cite{Goldman2000EnhancingSL}.

\subsection{Co-regularización}

Una ultima idea, dentro de los métodos envolventes basados en desacuerdo, pero algo alejada del concepto del psuodoetiquetado, es la de \textbf{co-regularización}. Al igual que en los métodos anteriores se entrenan diferentes clasificadores, sin embargo aqui el pseoudoetiquetado no es un proceso iterativo, si no que se introduce un término de regularización en la función de pérdida de cada clasificador que penaliza el desacuerdo entre los clasificadores en sus predicciones para los datos no etiquetados, de esta manera se fuerza a los clasificadores a estar de acuerdo en sus predicciones para los datos no etiquetados, lo que puede mejorar el rendimiento del modelo \cite{sindhwani2008rkhs}.

Veamos que quiere decir esto, para ello tenemos que definir el concepto de función de pérdida. En el aprendizaje supervisado, la función de pérdida es una medida de la discrepancia entre las predicciones del modelo y las etiquetas reales de los datos etiquetados, y se utiliza para guiar el proceso de entrenamiento del modelo:

\begin{defi}
    Bajo el marco del aprendizaje supervisado, sea $\mathbf{x} \in \mathcal{X}$ un vector de características, $y \in \mathcal{Y}$ su etiqueta correspondiente, y $f: \mathcal{X} \rightarrow \mathcal{Y}$ un modelo supervisado, se define la función de pérdida $L: \mathcal{Y} \times \mathcal{Y} \rightarrow \mathbb{R}$ como una función que mide la discrepancia entre la predicción del modelo $f(\mathbf{x})$ y la etiqueta real $y$, es decir, $L(f(\mathbf{x}), y)$ representa el costo o penalización por predecir $f(\mathbf{x})$ cuando la etiqueta real es $y$.
\end{defi}

Una posible idea que surge a partir de la función de pérdida es la de buscar un modelo que minimice la función de pérdida en el conjunto de datos etiquetados, es decir, encontrar un modelo $f$ que minimice la siguiente expresión:

\begin{equation}
    \sum_{i=1}^{l} L(f(\mathbf{x}_i), y_i)
\end{equation}

Sin embargo, estos modelos tienden a sobreajustar los datos etiquetados, perdiendo asi capacidad de generalización, para evitar esto surge el marco de la \textit{regularización} en el cual se introduce un término de regularización en la función de pérdida que penaliza la complejidad del modelo, de esta manera se busca un modelo que minimice la siguiente expresión:

\begin{equation}
    \sum_{i=1}^{l} L(f(\mathbf{x}_i), y_i) + \lambda \Omega(f)
\end{equation}

Dentro de este marco de regularización, se encuentra la idea de co-regularización, donde se introduce un término de regularización compuesto por la suma de un término de pérdida para cada clasificador supervisado y un término de regularización que penaliza el desacuerdo entre los clasificadores en sus predicciones para los datos no etiquetados, es decir, se busca un conjunto de clasificadores $f_1, f_2, \ldots, f_k$ que minimice la siguiente expresión:

\begin{equation}
    \sum_{u=1}^{k} \left(\sum_{i=1}^{l} L(f_u(\mathbf{x}_i), y_i) + \lambda \Omega(f_u) \right) + \gamma \sum_{u,v=1}^{k} \sum_{i=l+1}^{l+u} L(f_u(\mathbf{x}_i), f_v(\mathbf{x}_i))
\end{equation}

Este último término del sumatorio representa ese pseudoetiquetado, las etiquetas generadas por cada clasificador para los datos no etiquetados se utilizan para entrenar a los otros clasificadores en cada iteración, todo mientras se busca minimizar el desacuerdo entre los clasificadores en sus predicciones para los datos no etiquetados.

Destacar que aunque la co-regularización no se basa en un proceso iterativo de pseudoetiquetado, y aunque haya aparecido una función de pérdida que se minimiza, sigue siendo un método envolvente basado en desacuerdo, pues se siguen entrenando diferentes clasificadores supervisados con la información de los datos no etiquetados, y se sigue buscando aprovechar las diferencias entre los clasificadores para mejorar el rendimiento del modelo.

\section{Boosting}

Una última familia de métodos envolventes son aquellos que se basan en la idea del boosting. El boosting es una técnica de aprendizaje supervisado que se encuentra en el marco de los métodos de ensamblado, los cuales se basan en la idea de entrenar múltiples modelos supervisados y luego combinar sus predicciones para obtener una predicción final. El boosting se diferencia de otros métodos de ensamblado en que los modelos supervisados se entrenan de manera secuencial, cada modelo se entrena para corregir los errores del modelo anterior, lo que puede mejorar el rendimiento del modelo \cite{zhou2025ensemble}.

En el contexto del aprendizaje semi-supervisado se han desarrollado varios métodos basados en el boosting, pues permiten introducir una etapa de pseudoetiquetado tras cada entrenamiento. El precursos de estos modelos fue SSMBoost \cite{grandvalet2001boosting}, sin embargo, este método se basa en usar un algoritmo semi-supervisado de base, no en introducir el pseudoetiquetado tras cada entrenamiento, por lo que al no actuar de caja negra no es realmente un método envolvente. Posteriormente, se han desarrollado otros métodos basados en el boosting que si introducen el pseudoetiquetado tras cada entrenamiento, como por ejemplo Asssemble \cite{bennett2002exploiting} o SemiBoost \cite{mallapragada2008semiboost}, ambos se diferencian en la forma de introducir el pseudoetiquetado en la siguiente iteración. Assemble toma un enfoque más directo pues hace una selección aleatoria con remplazamiento (\textit{bootstraping}) sobre el conjunto formado por los datos etiquetados y pseudoetiquetados para generar el conjunto de datos etiquetados para la siguiente iteración. Por otro lado en SemiBoost se seleccionan las predicciones de mayor confianza para generar el conjunto de datos etiquetados para la siguiente iteración, la confianza de cada predicción individual se calcula usando un grafo local construido mediante similitud entre los datos de entrada.


\chapter{Preprocesado no-supervisado de datos}

La siguiente gran familia que vamos a presentar es la de los métodos de preprocesado no-supervisado de datos. La idea de esta familia es usar los datos no etiquetados en un paso previo al entrenamiento del modelo supervisado, con el objetivo de mejorar el rendimiento del modelo o reducir su complejidad \cite{van2020survey}.

Esta idea contrasta tanto con los métodos envolvenentes, que introducen los datos no etiquetados en el entrenamiento mediante el pseudoetiquetado, como con los métodos de intrínsecos, que introducen los datos no etiquetados en el entrenamiento mediante la adición de términos no etiquetados a la función de pérdida. Esta diferencia es importante, ya que permite a los métodos de preprocesado no-supervisado de datos ser compatibles con cualquier modelo supervisado, sin necesidad de modificar su función de pérdida o su proceso de entrenamiento.

Cada método de esta familia se distingue del resto por como implementa este \textit{paso previo}, aun así podemos identificar tres grandes categorías o subfamilias: Las basadas en la extracción de características, las basadas en el clustering y las basadas en el pre-entrenamiento de modelos.

\section{Extracción de características}

Los métodos de extracción de características buscan transformar los datos de entrada de su representación original a una representación diferente que mejore el rendimiento del modelo o reduzca la complejidad del entrenamiento. Esta transformación se puede realizar de diferentes maneras, como por ejemplo reduciendo la dimensionalidad de los datos, seleccionando características o generando nuevas características a partir de las originales. 

Uno de los métodos más conocidos de esta categoría es el \textbf{Análisis de Componentes Principales} (PCA), que busca encontrar una nueva representación de los datos en un espacio de menor dimensionalidad, donde las nuevas dimensiones son combinaciones lineales de las dimensiones originales que capturan la mayor parte de la varianza de los datos \cite{abdi2010principal}. El porqué de usar la varianza como criterio viene del siguiente teorema \cite{tour2006geometric}:

\begin{theo}
    Sea $\mathcal{X} \subseteq \mathbb{R}^d$ el dominio de las instancias y $\mathcal{D} = \{\mathbf{x}_1, \mathbf{x}_2, \ldots, \mathbf{x}_n\}$ un conjunto de datos. Entonces, la proyección de los datos sobre un subespacio de dimensión $d' < d$ que maximiza la varianza de los datos proyectados es la misma proyección que minimiza la distancia entre los datos originales y su aproximación en el nuevo subespacio, es decir, la que minimiza la siguiente expresión: 
    $$E = \frac{1}{n} \sum_{i=1}^n \|\mathbf{x}_i - \tilde{\mathbf{x}}_i\|^2$$ 
    donde $\tilde{\mathbf{x}}_i$ es la proyección ortogonal de $\mathbf{x}_i$ sobre el subespacio de dimensión $d'$.
\end{theo}

\begin{proof}
    \linebreak
    Sin perder generalidad, podemos asumir que los datos están centrados en el origen, es decir, que $\frac{1}{n} \sum_{i=1}^n \mathbf{x}_i = \mathbf{0}$.

    Sea $\{\mathbf{v}_j\}_{j=1}^{d'}$ una base ortonormal arbitraria que genera este nuevo subespacio. La proyección ortogonal de una muestra original $\mathbf{x}_i$ sobre dicho subespacio viene dada por la siguiente combinación lineal: 
    $$\tilde{\mathbf{x}}_i = \sum_{j=1}^{d'} (\mathbf{x}_i \cdot \mathbf{v}_j) \mathbf{v}_j$$

    Cualquier punto original puede descomponerse en la suma de su proyección y su diferencia o residuo: $\mathbf{x}_i = \tilde{\mathbf{x}}_i + (\mathbf{x}_i - \tilde{\mathbf{x}}_i)$. Si calculamos la norma al cuadrado de este vector desarrollando el producto escalar, obtenemos:
    $$\|\mathbf{x}_i\|^2 = \|\tilde{\mathbf{x}}_i\|^2 + \|\mathbf{x}_i - \tilde{\mathbf{x}}_i\|^2 + 2\big(\tilde{\mathbf{x}}_i \cdot (\mathbf{x}_i - \tilde{\mathbf{x}}_i)\big)$$

    Dado que $\tilde{\mathbf{x}}_i$ es la proyección ortogonal sobre el subespacio, el vector diferencia $(\mathbf{x}_i - \tilde{\mathbf{x}}_i)$ es perpendicular a dicho subespacio. En consecuencia, su producto escalar se anula ($\tilde{\mathbf{x}}_i \cdot (\mathbf{x}_i - \tilde{\mathbf{x}}_i) = 0$), lo que nos permite aplicar directamente el Teorema de Pitágoras: 
    $$\|\mathbf{x}_i\|^2 = \|\tilde{\mathbf{x}}_i\|^2 + \|\mathbf{x}_i - \tilde{\mathbf{x}}_i\|^2$$

    Despejando la distancia al cuadrado $\|\mathbf{x}_i - \tilde{\mathbf{x}}_i\|^2$ y sustituyéndola en la expresión de nuestra función a minimizar $E$, obtenemos: 
    $$ \frac{1}{n} \sum_{i=1}^n \|\mathbf{x}_i - \tilde{\mathbf{x}}_i\|^2 = \frac{1}{n} \sum_{i=1}^n \|\mathbf{x}_i\|^2 - \frac{1}{n} \sum_{i=1}^n \|\tilde{\mathbf{x}}_i\|^2 $$

    El primer término del lado derecho de la ecuación depende únicamente de la magnitud de los datos originales, por lo que es una constante estrictamente independiente de la elección de la base $\{\mathbf{v}_j\}$. Por lo tanto, el subespacio que minimiza la distancia a los datos originales (lado izquierdo) es aquel que maximiza el segundo término. Dado que este segundo término representa exactamente la varianza de los datos proyectados, demostramos que el subespacio óptimo es el que maximiza dicha varianza.
\end{proof}

Quedando establecido el porque se toma este criterio, el PCA se implementa calculando la matriz de covarianza de los datos $C = \frac{1}{n} \sum_{i=1}^n \mathbf{x}_i \mathbf{x}_i^T$. Pues, maximizar la varianza es equivalente a maximizar la expresión $\sum_{j=1}^{d'} \mathbf{v}_j^T C \mathbf{v}_j$ bajo la restricción de que los vectores $\{\mathbf{v}_j\}$ sean ortonormales. Este problema de optimización se resuelve usando métodos analíticos y su solución son los autovectores de la matriz de covarianza $C$, de los cuales se eligen los $d'$ de mayor valor propio \cite{tour2006geometric}.

El PCA es solo un ejemplo de método de extracción de características, pero existen muchos otros, como los métodos basados en autoencoders, que buscan aprender una representación comprimida de los datos mediante redes neuronales \cite{LI2023110176}. Lo que todos tienen en común es la idea base de hacer uso de los datos no etiquetados para transformar la representación de los datos de entrada, con el objetivo de mejorar el rendimiento del modelo supervisado o reducir su complejidad. 

\section{Cluster-then-label}

Los métodos de cluster-then-label cambian la idea anterior de realizar una transformación sobre los datos, por la idea de realizar una partición de los datos en clusters, y utilizar esta división del espacio para obtener mejoras en el rendimiento del modelo supervisado. Es decir, el \textit{paso previo} en esta categoría consiste en aplicar un algoritmo de clustering sobre todo el conjunto de datos, independientemente de si estan etiquetados o no \cite{van2020survey}. 

Desde una perspectiva formal, este enfoque transforma el problema de encontrar una única función de clasificación compleja en el problema de encontrar múltiples funciones más simples sobre subconjuntos del espacio. Sea $\mathcal{X} \subseteq \mathbb{R}^d$ el espacio de características y $\mathcal{Y}$ el espacio de etiquetas. Dado el conjunto de datos total $D = D_l \cup D_u$, el proceso se divide matemáticamente en dos fases:

\begin{enumerate}
    \item \textbf{Fase de partición (Clustering):} Se aplica un algoritmo de aprendizaje no supervisado sobre $D$ para encontrar la partición del espacio $\mathcal{X}$ en clústeres $\{C_1, C_2, \dots, C_K\}$. Como resultado de este proceso, se obtiene una función de asignación $f_c: \mathcal{X} \rightarrow \{1, 2, \dots, K\}$ que mapea cada punto de manera unívoca a su clúster correspondiente, de tal forma que $C_k = \{x \in \mathcal{X} \mid f_c(x) = k\}$.

    \item \textbf{Fase de clasificación (Labeling):} En lugar de optimizar una hipótesis global $h: \mathcal{X} \rightarrow \mathcal{Y}$ en un espacio de hipótesis $\mathcal{H}$ excesivamente complejo, se restringe el problema. Para cada clúster $C_k$, se entrena una hipótesis local $h_k \in \mathcal{H}_k$ utilizando únicamente el subconjunto de datos etiquetados que caen dentro de esa región, es decir, el conjunto $\{(x,y) \in D_l \mid f_c(x) = k\}$.
\end{enumerate}

Finalmente, el clasificador global semi-supervisado se define como una función definida a trozos (\textit{piecewise function}), donde la predicción para una nueva instancia $x$ viene dada por el modelo local de su clúster correspondiente:

\begin{equation}
    h(x) = \sum_{k=1}^K h_k(x) \cdot \mathbb{I}(f_c(x) = k)
\end{equation}

donde $\mathbb{I}(\cdot)$ es la función indicadora. Esta formulación permite que la estructura subyacente revelada por los datos no etiquetados (la partición) reduzca drásticamente la complejidad funcional del aprendizaje supervisado posterior.

Al igual que antes existen diferentes enfoques a la hora de aplicar esta idea que varian en función del algoritmo de clustering elegido y de la forma en la que se utiliza la partición obtenida para mejorar el rendimiento del modelo supervisado.

Un posible enfoque es el de \textit{Multi-Manifold Semi-Supervised Learning, Goldberg et al., 2009} \cite{pmlr-v5-goldberg09a}, que se basa en la idea de que los datos pueden estar distribuidos en diferentes subespacios o variedades, y que cada una de estas variedades puede ser modelada por un algoritmo de clustering diferente. En este enfoque, se aplica un algoritmo de clustering no supervisado para identificar estas variedades, y luego se entrena un modelo supervisado para cada variedad utilizando solo las muestras etiquetadas que pertenecen a dicha variedad. De esta forma, se obtiene un conjunto de modelos supervisados especializados en cada variedad, lo que puede mejorar el rendimiento del modelo global.

O el de \textit{Semi-Supervised Clustering Using Genetic Algorithms, Demiriz et al., 1999} \cite{demiriz1999semi} donde se utuliza un algoritmo de clustering semi-supervisado que combina tanto la información de los datos no etiquetados como la información de las muestras etiquetadas para realizar la partición del espacio.

En general, el concepto siempre es el mismo, hacer uso de los datos no etiquetados para realizar una partición del espacio, y luego utilizar esta partición para mejorar el rendimiento del modelo supervisado. Sin embargo, las formas de implementar esta idea pueden variar mucho, y cada una de ellas puede ser más adecuada para ciertos tipos de datos o problemas específicos.

\section{Pre-entrenamiento de modelos}

Los métodos de pre-entrenamiento de modelos buscan aprovechar los datos no etiquetados para obtener una mejor inicialización de los parámetros del modelo supervisado que se va a entrenar. La idea es entrenar un modelo no supervisado utilizando los datos no etiquetados, y luego utilizar los parámetros aprendidos por este modelo como punto de partida para el entrenamiento del modelo supervisado \cite{van2020survey}.

Este tipo de métodos es especialmente útil en el contexto de las redes neuronales profundas, donde la inicialización de los parámetros puede tener un gran impacto en el rendimiento del modelo supervisado, pues los algoritmos de entrenamiento pueden converger lentamenete o presentar problemas de optimización. La idea general consiste en entrenar los parámetros de las capas inferiores de la red neuronal profunda utilizando los datos no etiquetados, y luego fine-tunear el modelo con los datos etiquetados \cite{hinton2006fast,vincent2008extracting}.

Para ilustrar este proceso desde un punto de vista matemático, el enfoque más representativo es el uso de \textbf{Autoencoders} (o codificadores automáticos) apilados. Un autoencoder es una red neuronal diseñada para aprender la función identidad, es decir, para reconstruir su propia entrada, pero forzando al modelo a aprender una representación latente útil \cite{van2020survey}.

Formalmente, dado un conjunto de datos no etiquetados $D_u = \{x_1, \dots, x_m\}$ donde $x_i \in \mathbb{R}^d$, un autoencoder estándar se compone de dos funciones \cite{LI2023110176}:

\begin{enumerate}
    \item \textbf{El codificador (Encoder) $h$}: Mapea el vector de entrada $x$ a una representación latente oculta $z \in \mathbb{R}^{d'}$, típicamente de menor dimensión ($d' < d$). Se define como:
    \begin{equation}
        z = h(x) = \sigma(W_1 x + b_1)
    \end{equation}
    donde $W_1 \in \mathbb{R}^{d' \times d}$ es la matriz de pesos, $b_1 \in \mathbb{R}^{d'}$ es el vector sesgo, y $\sigma$ es una función de activación no lineal (como sigmoide o ReLU).
    
    \item \textbf{El decodificador (Decoder) $g$}: Intenta reconstruir la entrada original a partir de la representación latente. Se define como:
    \begin{equation}
        \hat{x} = g(z) = \sigma(W_2 z + b_2)
    \end{equation}
    donde $W_2 \in \mathbb{R}^{d \times d'}$ y $b_2 \in \mathbb{R}^d$.
\end{enumerate}

El pre-entrenamiento no supervisado consiste en encontrar los parámetros $\theta = \{W_1, b_1, W_2, b_2\}$ que minimizan una función de pérdida que penaliza el error de reconstrucción. La función de pérdida más común es el Error Cuadrático Medio (MSE):

\begin{equation}
    \mathcal{L}_{AE}(\theta) = \frac{1}{m} \sum_{i=1}^{m} ||x_i - g(h(x_i))||^2
\end{equation}

Al minimizar esta función utilizando los datos no etiquetados, los pesos $W_1$ y $b_1$ del codificador aprenden a capturar la estructura subyacente de los datos. En el contexto del pre-entrenamiento de redes profundas, una vez que este autoencoder es entrenado, se descarta el decodificador $g$, y los parámetros optimizados $W_1$ y $b_1$ se utilizan como la \textbf{inicialización exacta} para la primera capa oculta del modelo supervisado final. Este proceso puede repetirse secuencialmente capa por capa (Autoencoders apilados) antes de aplicar el ajuste fino (\textit{fine-tuning}) con los datos etiquetados mediante retropropagación estándar \cite{vincent2008extracting}.

Desde el punto de vista teórico las difentes capas de las redes neuronales profundas pueden ser vistas como diferentes niveles de abstracción en la representación de los datos, donde las capas inferiores capturan características más simples y las capas superiores capturan características más complejas. Por lo tanto, el pre-entrenamiento de modelos busca empujar al modelo hacia una mejor representación de los datos desde el principio del entrenamiento. Por lo tanto estos metodos estan intimamente relacionados con los métodos de extracción de características, aunque con la diferencia de que el pre-entrenamiento de modelos se centra en la inicialización de los parámetros del modelo supervisado que son posteriomente modificados en el entrenamietno, mientras que en los métodos de extracción de características la transformación de la representación de los datos de entrada realizada es permanente \cite{van2020survey}. 

\chapter{Métodos intrínsecamente semi-supervisados}

En este capítulo se describen los métodos intrínsecamente semi-supervisados, es decir, aquellos que desde su concepción están diseñados para entrenar de manera conjunta con datos etiquetados y no etiquetados. Su rasgo distintivo es que los ejemplos no etiquetados no aparecen como una etapa auxiliar o previa, sino que se incorporan de forma explícita en la formulación del método: bien a través de un modelo probabilístico (por ejemplo, en la verosimilitud de $p(\mathbf{x},y)$), o bien como términos adicionales en la función objetivo (por ejemplo, regularización, consistencia o suavidad).

Siguiendo la última división de los métodos inductivos presentada en la taxonomía del Capítulo 3, organizamos los métodos intrínsecos en tres grandes grupos: métodos generativos, métodos discriminativos y métodos basados en variedades. En cada caso se introduce la intuición del enfoque y se presentan formulaciones representativas: mezclas de modelos optimizadas con EM, clasificadores de máximo margen como las S3VM y métodos basados en perturbaciones como las \textit{ladder networks}, y, finalmente, técnicas de regularización de la variedad mediante grafos.

\section{Generativos}

Los métodos generativos constituyen uno de los enfoques clásicos del aprendizaje semi-supervisado \cite{zhu2005semi}. Su idea central es modelar explícitamente el proceso que genera los datos, esto es, la distribución conjunta $p(\mathbf{x},y)$, y, una vez estimada, usarla de diferentes formas para clasificar \cite{van2020survey}. Este enfoque conecta directamente con la discusión del Capítulo 2 sobre cómo la información de $p(\mathbf{x})$ puede ayudar a inferir $p(y\mid\mathbf{x})$ (Sección \ref{sec:suposiciones}).

Dentro de esta familia, el referente clásico son los modelos de mezcla optimizados mediante máxima verosimilitud, normalmente con el algoritmo EM cuando hay variables latentes \cite{dempster1977maximum}.

\subsection{Mezcla de modelos}

Los modelos de mezcla asumen que los datos se generan a partir de una combinación de componentes probabilísticas, típicamente una por clase, con parámetros desconocidos \cite{van2020survey}. Bajo este supuesto, los datos etiquetados y no etiquetados pueden contribuir conjuntamente a estimar el modelo.

La intuición en semi-supervisado es que los datos no etiquetados ayudan a identificar la estructura de la mezcla (componentes, pesos y forma), mientras que los pocos datos etiquetados permiten asociar componentes a clases \cite{zhu2005semi,zhu2009introduction}.

\subsubsection{Mezcla de modelos supervisada}

Dada una instancia $\mathbf{x}$, la predicción se obtiene tomando la clase más probable a posteriori:

\begin{equation}
    \hat{y} = \operatorname*{arg\,max}_{y} p(y|\mathbf{x})
\end{equation}

Para realizar dicho cálculo, se emplea la regla de Bayes:

\begin{equation}
    p(y|\mathbf{x}) = \frac{p(\mathbf{x}|y)p(y)}{\sum_{y'}p(\mathbf{x}|y')p(y')}
\end{equation}

por lo que el problema se reduce a estimar $p(\mathbf{x}|y)$ y $p(y)$. En el contexto de una mezcla de componentes probabilísticos, esto equivale a estimar los parámetros que definen dichas componentes.  

Si denotamos por $\theta$ al vector con los parámetros que definen los modelos y por $D_l=\{(\mathbf{x}_i,y_i)\}_{i=1}^{l}$ al conjunto etiquetado. Entonces el estimador de máxima verosimilitud de los parámetros es

\begin{equation}
    \hat{\theta} = \operatorname*{arg\,max}_{\theta} p(D_l|\theta) = \operatorname*{arg\,max}_{\theta} \log p(D_l|\theta).
\end{equation}

Es decir, buscamos los parámetros $\theta$ que maximizan la verosimilitud de los datos de entrenamiento. Además, como estamos tratando con datos independientes e idénticamente distribuidos, la verosimilitud se factoriza como el producto de las probabilidades individuales, lo que nos lleva a la siguiente expresión para la log-verosimilitud:

\begin{equation}
    \log p(D_l|\theta) = \log \prod_{i=1}^{l} p(\mathbf{x}_i, y_i|\theta)
     = \sum_{i=1}^{l} \log \bigl(p(y_i|\theta)p(\mathbf{x}_i|y_i,\theta)\bigr).
\end{equation}

Para familias concretas (por ejemplo, mezclas gaussianas con etiquetas observadas), este problema suele admitir estimadores cerrados o de optimización relativamente directa \cite{zhu2009introduction}.

\subsubsection{Mezcla de modelos semi-supervisada}

En el caso semi-supervisado se dispone, además de $D_l$, de un conjunto no etiquetado $D_u=\{\mathbf{x}_j\}_{j=l+1}^{l+u}$. El conjunto completo es $D=D_l\cup D_u$, y la log-verosimilitud pasa a ser:

\begin{equation}
    \log p(D|\theta) = \sum_{i=1}^{l} \log p(\mathbf{x}_i,y_i|\theta) + \sum_{j=l+1}^{l+u} \log p(\mathbf{x}_j|\theta).
\end{equation}

Donde en el término de los datos no etiquetados se incorpora la distribución marginal
\begin{equation}
    p(\mathbf{x}_j|\theta)=\sum_{y\in\mathcal{Y}} p(\mathbf{x}_j,y|\theta)=\sum_{y\in\mathcal{Y}} p(y|\theta)p(\mathbf{x}_j|y,\theta),
\end{equation}
que refleja que la etiqueta de $\mathbf{x}_j$ es desconocida \cite{zhu2009introduction}.

En general, esta optimización deja de ser cóncava y no suele resolverse de forma analítica. Por ello, el método estándar es EM, un procedimiento iterativo de máxima verosimilitud con variables latentes \cite{dempster1977maximum}. De forma intuitiva, EM repite el siguiente ciclo:
\begin{itemize}
    \item \textbf{Inicialización}: partir de unos parámetros iniciales $\theta^{(0)}$ (por ejemplo, obtenidos con los datos etiquetados).
    \item \textbf{E-step}: con los parámetros actuales, estimar para cada dato no etiquetado la probabilidad de pertenecer a cada clase, es decir, una etiqueta \textit{blanda}.
    \item \textbf{M-step}: de manera análoga al caso supervisado, reestimar los parámetros del modelo, pero usando ahora las etiquetas probabilísticas obtenidas en la E-step como pesos en la máxima verosimilitud.
    \item \textbf{Iteración}: repetir E-step y M-step hasta convergencia.
\end{itemize}

Esta dinámica es conceptualmente parecida al pseudoetiquetado de los métodos envolventes (Capítulo \ref{cap:envolventes}) \cite{zhu2009introduction,zhou2010semi}, pero aquí las asignaciones para $D_u$ son probabilísticas, mientras que en los métodos envolventes se termina obteniendo una etiqueta final para cada ejemplo.

La mezcla de modelos es un enfoque que puede ser muy eficaz cuando el modelo generativo está bien especificado. Sin embargo, su rendimiento depende críticamente de hipótesis como la identificabilidad de la mezcla y la corrección del modelo; si estas fallan, los datos no etiquetados pueden degradar el rendimiento \cite{zhu2005semi,van2020survey}.


\section{Discriminativos}

Esta familia de métodos se distingue de la anterior en que no trata de modelar el proceso generativo ni la densidad conjunta $p(\mathbf{x},y)$, sino que se centra directamente en aprender un modelo discriminativo $p(y\mid\mathbf{x})$ o una función de decisión $f:\mathcal{X}\to\mathcal{Y}$ \cite{ng2001discriminative}. En este caso, la información de los datos no etiquetados se incorpora con el objetivo de mejorar la generalización del modelo, normalmente a través de suposiciones adicionales sobre la distribución subyacente de los datos (por ejemplo, la separación en baja densidad o la suposición de suavidad) \cite{van2020survey}.
\subsection{Máximo margen}
\label{subsec:maximo_margen}

Los métodos de máximo margen buscan aprender una frontera de decisión que quede lo más alejada posible de los ejemplos de entrenamiento. En el caso de clasificación binaria, esto suele formularse como la maximización del margen geométrico del clasificador \cite{van2020survey}. En aprendizaje semi-supervisado, esta idea se combina con la suposición de separación en baja densidad: además de clasificar bien los ejemplos etiquetados, se prefiere que la frontera atraviese regiones del espacio de entrada donde haya pocos datos \cite{van2020survey,pmlr-v5-goldberg09a}.

El representante más conocido es la \textit{Semi-Supervised Support Vector Machine} (S3VM), también denominada históricamente \textit{Transductive SVM} (TSVM) \cite{10.7551/mitpress/9780262033589.001.0001}. Para introducirla, repasamos primero la formulación básica de las SVM.

\subsubsection{SVM} 

Consideramos un conjunto etiquetado $\{(\mathbf{x}_i,y_i)\}_{i=1}^l$ con $\mathbf{x}_i\in\mathbb{R}^d$ e $y_i\in\{-1,1\}$, y un clasificador lineal de la forma
\begin{equation}
f(\mathbf{x})=\mathbf{w}^{\top}\mathbf{x}+b,
\end{equation}
donde $\mathbf{w}\in\mathbb{R}^d$ y $b\in\mathbb{R}$ \cite{zhu2009introduction}. Por lo que la frontera de decisión es el siguiente hiperplano:
\begin{equation}
\left\{\mathbf{x}\in\mathbb{R}^d: \mathbf{w}^{\top}\mathbf{x}+b=0\right\},
\end{equation}
mientras que la predicción se obtiene con $\mathrm{sign}(f(\mathbf{x}))$.

La distancia (no firmada) de un punto $\mathbf{x}$ a dicho hiperplano es $\frac{|f(\mathbf{x})|}{\|\mathbf{w}\|}$, y para los ejemplos etiquetados se toma la distancia con signo que es $\frac{y_i f(\mathbf{x}_i)}{\|\mathbf{w}\|}$. Con todo esto, el margen geométrico del clasificador se define como
\begin{equation}
\gamma = \min_{i=1,\dots,l} \frac{y_i\bigl(\mathbf{w}^{\top}\mathbf{x}_i+b\bigr)}{\|\mathbf{w}\|}.
\end{equation}
Primero, vamos a asumir que los datos son \textit{linealmente separables}, es decir, que existe un hiperplano tal que $y_i\bigl(\mathbf{w}^{\top}\mathbf{x}_i+b\bigr)>0$ para todo $i=1,\dots,l$. Bajo esta suposición, el \textit{hard-margin SVM} busca el hiperplano separador que maximiza $\gamma$, lo que puede escribirse como
\begin{equation}
\max_{\mathbf{w},b} \quad \min_{i=1,\dots,l} \frac{y_i\bigl(\mathbf{w}^{\top}\mathbf{x}_i+b\bigr)}{\|\mathbf{w}\|}.
\end{equation}

Este problema es invariante frente a escalado: si multiplicamos $(\mathbf{w},b)$ por una constante $c>0$, la frontera de decisión $\{\mathbf{x}: \mathbf{w}^{\top}\mathbf{x}+b=0\}$ no cambia. Por ello, podemos fijar una escala imponiendo $\min_i y_i\bigl(\mathbf{w}^{\top}\mathbf{x}_i+b\bigr)=1$, equivalente a exigir $y_i\bigl(\mathbf{w}^{\top}\mathbf{x}_i+b\bigr)\ge 1$ para todo $i$. En consecuencia, maximizar el margen es equivalente a maximizar $\frac{1}{\|\mathbf{w}\|}$, o de forma equivalente minimizar la norma de $\mathbf{w}$:
\begin{equation}
\begin{aligned}
\min_{\mathbf{w},b} \quad & \frac{1}{2}\|\mathbf{w}\|^2 \\
\text{s.a.} \quad & y_i\bigl(\mathbf{w}^{\top}\mathbf{x}_i+b\bigr) \geq 1, \quad i=1,\dots,l.
\end{aligned}
\end{equation}

Cuando los datos no son separables (esto es, no existe un hiperplano que separe perfectamente las clases), se introduce el \textit{soft margin} mediante variables de holgura $\xi_i\ge 0$ y un parámetro $C>0$ que controla el compromiso entre maximizar el margen y minimizar errores. Intuitivamente, $\xi_i$ mide cuánto viola el ejemplo $i$ la restricción de margen: $\xi_i=0$ significa que se cumple el margen, $0<\xi_i<1$ indica que el punto queda dentro del margen pero sigue estando bien clasificado, y $\xi_i>1$ corresponde a una clasificación errónea.
\begin{equation}
\begin{aligned}
\min_{\mathbf{w},b} \quad & \frac{1}{2}\|\mathbf{w}\|^2 + C \sum_{i=1}^l \xi_i \\
\text{s.a.} \quad & y_i\bigl(\mathbf{w}^{\top}\mathbf{x}_i+b\bigr) \geq 1 - \xi_i, \quad i=1,\dots,l, \\
                  & \xi_i \geq 0.
\end{aligned}
\end{equation}

Finalmente, puede expresarse en el marco de la minimización de riesgo regularizado (Sección \ref{sec:perdida_regularizacion}). Despejando $\xi_i$ de $y_i f(\mathbf{x}_i) \ge 1-\xi_i$ obtenemos $\xi_i\ge 1-y_if(\mathbf{x}_i)$, y junto con $\xi_i\ge 0$ se tiene $\xi_i=\max\bigl(0,1-y_if(\mathbf{x}_i)\bigr)$. Por lo tanto, la función objetivo se reescribe como
\begin{equation}
\min_{\mathbf{w},b} \quad \frac{1}{2}\|\mathbf{w}\|^2 + C \sum_{i=1}^l \max\bigl(0, 1 - y_i f(\mathbf{x}_i)\bigr).
\end{equation}
En este caso la función de pérdida es $c(\mathbf{x},y,f(\mathbf{x})) = \max\bigl(0, 1 - y f(\mathbf{x})\bigr)$ (pérdida \textit{hinge}) y el regularizador es $\frac{1}{2}\|\mathbf{w}\|^2$.

\subsubsection{S3VM}

La idea central de las S3VM es incorporar un conjunto no etiquetado $\{\mathbf{x}_j\}_{j=l+1}^{l+u}$ penalizando que dichos puntos queden cerca de la frontera: se prefiere que $|f(\mathbf{x}_j)|$ sea grande, de modo que la frontera se sitúe en regiones de baja densidad \cite{van2020survey,10.7551/mitpress/9780262033589.001.0001}. Una forma estándar de derivar esta penalización es introducir etiquetas latentes $\hat{y}_j\in\{-1,1\}$ para los ejemplos no etiquetados y añadir un término de pérdida \textit{hinge} $\max\bigl(0,1-\hat{y}_j f(\mathbf{x}_j)\bigr)$, que al minimizar respecto a $\hat{y}_j$ se convierte en $\max\bigl(0,1-|f(\mathbf{x}_j)|\bigr)$, conocida como \textit{hat loss} \cite{10.7551/mitpress/9780262033589.001.0001}.

Con ello, una formulación típica de S3VM es
\begin{equation}
\min_{\mathbf{w},b} \quad \frac{1}{2}\|\mathbf{w}\|^2 + C \sum_{i=1}^l \max\bigl(0, 1 - y_i f(\mathbf{x}_i)\bigr) + C^* \sum_{j=l+1}^{l+u} \max\bigl(0, 1 - |f(\mathbf{x}_j)|\bigr).
\end{equation}

En términos del marco de la Sección \ref{sec:perdida_regularizacion}, el sumatorio sobre los datos no etiquetados puede interpretarse como un término adicional dependiente de $\{\mathbf{x}_j\}$, que incorpora explícitamente la suposición de separación en baja densidad.

En la práctica suele ser necesario imponer alguna restricción adicional (por ejemplo, sobre el balance de etiquetas asignadas a los datos no etiquetados) para evitar soluciones degeneradas \cite{van2020survey}. El problema resultante es no convexo, por lo que se recurre a métodos heurísticos y aproximados; además, si la suposición de separación en baja densidad no se cumple, el uso de la información no etiquetada puede incluso degradar el rendimiento \cite{van2020survey}.

\subsection{Basados en perturbaciones}
Por otra parte, los métodos basados en perturbaciones se fundamentan en la suposición de suavidad: puntos cercanos en el espacio de entrada deberían tener etiquetas similares \cite{10.7551/mitpress/9780262033589.001.0001}. Para aprovechar esta suposición, se introducen perturbaciones o ruido estocástico en los datos de entrenamiento y se añade un término de regularización por consistencia a la función de pérdida. Este término penaliza que el modelo produzca predicciones diferentes para las versiones original y perturbada de un mismo ejemplo \cite{van2020survey}.

Este enfoque se ha popularizado especialmente en el contexto de las redes neuronales profundas, debido a su capacidad para aprender representaciones ricas y a la facilidad algorítmica de incorporar perturbaciones en el proceso de entrenamiento \cite{van2020survey}. Un ejemplo pionero es el método \textit{ladder network} \cite{rasmus2015semi}, que combina una arquitectura de red neuronal con un término de pérdida que evalúa la discrepancia entre las activaciones de la red para ejemplos limpios y perturbados.

\subsubsection{Redes neuronales}

Antes de proceder a exponer el método de las \textit{ladder networks}, repasamos brevemente el marco general de las redes neuronales, más concretamente el marco de las redes neuronales feedforward. 

Una red neuronal es un sistema que transforma un vector de entrada $\mathbf{x}\in\mathbb{R}^d$ en un vector de salida $f(\mathbf{x})\in\mathbb{R}^k$ mediante la propagación de la entrada por una serie de capas conectadas secuencialmente \cite{van2020survey}.

Cada capa de la red se compone de un conjunto de perceptrones, un modelo matemático de neurona. Cada uno de estos perceptrones recibe como entrada un vector $\mathbf{z}\in\mathbb{R}^m$, o bien el de entrada o bien las salidas de la capa anterior, y produce una salida escalar $a\in\mathbb{R}$ mediante la aplicación de una función de activación $\sigma:\mathbb{R}\to\mathbb{R}$ a una combinación lineal de las entradas, es decir, $a=\sigma(\mathbf{w}^{\top}\mathbf{z}+b)$, donde $\mathbf{w}\in\mathbb{R}^m$ y $b\in\mathbb{R}$ son los parámetros del perceptrón \cite{bishop2006pattern}.

En el contexto de aprendizaje supervisado, el entrenamiento de la red neuronal consiste en ajustar los parámetros de los perceptrones, representados por una matriz $W$ que contiene tanto los pesos $\mathbf{w}$ como los sesgos $b$, para minimizar una función de pérdida \cite{van2020survey}.

\begin{equation}
    \mathcal{L}(W) = \frac{1}{l} \sum_{i=1}^l c(\mathbf{x}_i,y_i,f(\mathbf{x}_i;W))
\end{equation}

Los parámetros se optimizan iterativamente mediante la retropropagación del error, normalmente con un método de descenso de gradiente estocástico o alguna de sus variantes \cite{NIPS2014_f033ed80}.

\subsubsection{Ladder networks}

Una vez introducidos los conceptos básicos de las redes neuronales, es posible presentar el método de las \textit{ladder networks} \cite{rasmus2015semi}. Este enfoque extiende la arquitectura de una red neuronal feedforward estándar incorporando un término de pérdida no supervisado, basado en el principio de los autoencoders eliminadores de ruido.

En concreto, el modelo se estructura en un codificador y un decodificador. Durante el entrenamiento, cada dato de entrada realiza dos pasadas por el codificador: una pasada limpia y una pasada a la que se le inyecta ruido en todas las capas ocultas. Posteriormente, el decodificador intenta ir limpiando este ruido hacia atrás, reconstruyendo las activaciones limpias originales capa por capa a partir de las representaciones perturbadas. De este modo, el término de pérdida global se define como la suma del coste de clasificación supervisado y un coste de consistencia o reconstrucción calculado a lo largo de todas las capas de la red \cite{rasmus2015semi}.

La función de pérdida total a optimizar se puede formular de la siguiente manera:
\begin{equation}
    \mathcal{L}(W) = \frac{1}{l} \sum_{i=1}^l c(\mathbf{x}_i,y_i,f(\mathbf{x}_i;W)) + \sum_{i=1}^{l+u} \sum_{k=1}^K ReconstructionLoss\bigl(\hat{\mathbf{z}}_i^k, \mathbf{z}_i^k\bigr)
\end{equation}

Donde $\mathbf{z}_i^k$ es la activación oculta de la capa $k$ en el codificador para el ejemplo limpio $\mathbf{x}_i$, mientras que $\hat{\mathbf{z}}_i^k$ es la activación reconstruida por el decodificador a partir de la versión perturbada. Finalmente, $\text{ReconstructionLoss}$ es una función que penaliza la discrepancia entre ambas (como, por ejemplo, la norma $L_2$ al cuadrado) \cite{rasmus2015semi}.

A través de esta arquitectura, las \textit{ladder networks} consiguen aprovechar simultáneamente tanto la señal de los datos etiquetados como la regularización proporcionada por los datos no etiquetados\footnote{Esta simultaneidad es, de hecho, lo que hace a las \textit{ladder networks} modelos intrínsecos, a diferencia del preentrenamiento por capas visto en la Sección \ref{sec:pre-training}, donde se usaba un \textit{encoder-decoder} como fase inicial independiente seguida de la fase supervisada.}, forzando a la red a extraer características latentes útiles y robustas al ruido que, en última instancia, mejoran la generalización del modelo predictivo \cite{rasmus2015semi}.
\section{Basados en variedades}
\label{sec: variedades}

Finalmente, los métodos basados en variedades constituyen una categoría independiente en la presente taxonomía. Por un lado, no pueden considerarse generativos, puesto que no modelan explícitamente el proceso de generación de los datos. Por otro lado, aunque (al igual que los métodos discriminativos) se centran en aprender una función de decisión, su enfoque subyacente es diferente al de las técnicas de máximo margen o de perturbaciones.

Estos métodos se apoyan en la hipótesis de la variedad: se asume que los datos observados se concentran cerca de una variedad de baja dimensión inmersa en el espacio de entrada original \cite{10.7551/mitpress/9780262033589.001.0001}. En consecuencia, el papel de los datos no etiquetados no es solo ajustar una frontera global, sino capturar la geometría de dicha variedad y orientar la búsqueda hacia funciones que respeten su estructura intrínseca \cite{van2020survey}.

\subsection{Técnicas de regularización de la variedad}

Una forma habitual de aproximar la variedad subyacente es construir un grafo de similitud entre instancias \cite{van2020survey}. Dado un conjunto de entrenamiento formado por $\{(\mathbf{x}_i, y_i)\}_{i=1}^l$ instancias etiquetadas y $\{\mathbf{x}_j\}_{j=l+1}^{l+u}$ instancias sin etiquetar, se construye un grafo no dirigido $G=(V,E)$ donde cada nodo corresponde a una instancia (etiquetada o no)\footnote{Como consecuencia, los grafos resultantes suelen tener gran tamaño cuando se dispone de una gran cantidad de instancias no etiquetadas \cite{zhu2009introduction}.}, y a cada arista $(i,j)\in E$ se le asigna un peso $w_{ij}$ que refleja la similitud entre $\mathbf{x}_i$ y $\mathbf{x}_j$. Denotamos por $W=(w_{ij})$ la matriz de pesos. La intuición es que, si $w_{ij}$ es alto, entonces es más probable que ambas instancias compartan etiqueta \cite{zhu2009introduction}.

Hay varias formas de definir la vecindad y/o la similitud. Por ejemplo, puede construirse un grafo de $k$ vecinos más cercanos (k-NN) y asignar pesos gaussianos $w_{ij} = \exp\bigl(-\frac{\|\mathbf{x}_i - \mathbf{x}_j\|^2}{2\sigma^2}\bigr)$ a las aristas \cite{zhu2009introduction}.

Una vez construido el grafo, se busca una función de decisión $f:\mathcal{X}\to\mathcal{Y}$ que no solo clasifique correctamente los ejemplos etiquetados, sino que también sea suave respecto a la estructura del grafo; es decir, que produzca salidas similares en nodos conectados por aristas con pesos altos \cite{zhu2009introduction}.

Para ello, un enfoque representativo es el presentado en \textit{On Manifold Regularization, Belkin et al., 2005} \cite{belkin2005manifold}. Partiendo de un marco de aprendizaje supervisado, se considera un kernel $K$\footnote{Un kernel es una función $K:\mathcal{X}\times\mathcal{X}\to\mathbb{R}$ que puede interpretarse como un producto escalar en un espacio de características (posiblemente de mayor dimensión) \cite{bishop2006pattern}.}, con un espacio de funciones asociado $\mathcal{H}_K$ y norma $\|\cdot\|_K$. Se busca una función $f\in\mathcal{H}_K$ que minimice la siguiente función de pérdida regularizada:
\begin{equation}
    \mathcal{L}(f) = \frac{1}{l} \sum_{i=1}^l c(\mathbf{x}_i,y_i,f(\mathbf{x}_i)) + \lambda \|f\|_K^2,
\end{equation}

donde el primer término es la pérdida de clasificación sobre los datos etiquetados y el segundo término es un regularizador que penaliza la complejidad de la función $f$ \cite{belkin2005manifold}.

Finalmente, la información no etiquetada se introduce mediante un término de regularización adicional, utilizando la matriz de pesos $W$, que penaliza la falta de suavidad de $f$ respecto a la estructura del grafo:
\begin{equation}
    \|f\|_I^2 = \frac{1}{2} \sum_{i,j} w_{ij} \bigl(f(\mathbf{x}_i) - f(\mathbf{x}_j)\bigr)^2,
\end{equation}

Este término puede interpretarse como una medida de la variación de $f$ a lo largo de las aristas del grafo, ponderada por los pesos de similitud. Además, si denotamos $n=l+u$ y definimos el vector $\mathbf{f} = \bigl(f(\mathbf{x}_1),\dots,f(\mathbf{x}_n)\bigr)^{\top}$, la matriz diagonal de grados $D$ con $D_{ii}=\sum_j w_{ij}$ y el laplaciano no normalizado $L=D-W$, entonces se obtiene una expresión matricial compacta para él mismo $\|f\|_I^2 = \mathbf{f}^{\top} L \mathbf{f}$ \cite{belkin2005manifold}.

De este modo, la función de pérdida total a minimizar es

\begin{equation}
    \mathcal{L}(f) = \frac{1}{l} \sum_{i=1}^l c(\mathbf{x}_i,y_i,f(\mathbf{x}_i)) + \lambda_A \|f\|_K^2 + \lambda_I \|f\|_I^2,
\end{equation}

donde $\lambda_A$ y $\lambda_I$ son hiperparámetros que controlan el peso relativo de cada término de regularización \cite{belkin2005manifold}.

Dependiendo de la elección del kernel $K$ y de la función de pérdida $c$, esta formulación puede dar lugar a diferentes algoritmos concretos. Un ejemplo clásico son las \textit{Laplacian SVM}, que combinan una pérdida \textit{hinge} con el término de regularización basado en el grafo. De este modo, a diferencia de las S3VM vistas en la sección anterior, la información no etiquetada se incorpora a través de la geometría del grafo y no mediante la penalización directa de puntos cercanos a la frontera de decisión \cite{van2020survey}.

En la práctica, estos métodos pueden ser costosos computacionalmente, ya que suelen requerir operar con matrices de tamaño $n\times n$ (donde $n=l+u$). En particular, en algunos casos el coste puede alcanzar $O(n^3)$, o bien $O(c\,n^2)$ para un $c \ll n$ \cite{van2020survey}. Además, el rendimiento es muy sensible a cómo se construye el grafo y a la elección y escala de los pesos $w_{ij}$; si el grafo no refleja bien la geometría de los datos, la información no etiquetada puede propagarse de forma perjudicial \cite{zhu2009introduction}.

\chapter{Métodos Transductivos}

Hasta ahora, las familias de métodos presentadas en los capítulos anteriores eran \emph{inductivas}: su objetivo era construir un modelo a partir de los datos de entrenamiento y, posteriormente, usarlo para predecir la clase de nuevas instancias del espacio de entrada. En este capítulo se presentan métodos \emph{transductivos}, que en lugar de aprender una regla de decisión general, se centran en inferir la clase de las instancias no etiquetadas disponibles en el propio conjunto de entrenamiento \cite{van2020survey}.

El enfoque transductivo aparece de manera natural en los métodos semi-supervisados basados en grafos. En ellos se construye un grafo de similitud donde cada nodo representa una instancia (etiquetada o no) y las aristas ponderadas codifican la proximidad entre pares de ejemplos. Estos métodos suelen asumir \emph{suavidad de etiquetas} sobre el grafo: si dos nodos están conectados por una arista con peso grande, entonces es razonable esperar que sus etiquetas (o puntuaciones de clase) sean similares \cite{zhu2005semi}.

Esta idea es conceptualmente cercana a la de los métodos basados en variedades (véase la Sección \ref{sec: variedades}), donde el grafo se utiliza como una aproximación discreta de la geometría de los datos. La diferencia clave es que, en el planteamiento transductivo, la función de decisión se define \emph{directamente sobre los nodos} del grafo, de modo que el problema se reduce a propagar de forma coherente la información de las etiquetas conocidas hacia los nodos no etiquetados \cite{van2020survey}.

Usando el marco de pérdida y regularización introducido en la Sección \ref{sec:perdida_regularizacion}, el aprendizaje transductivo sobre grafos puede formularse como la minimización de un término de ajuste en los nodos etiquetados más un término que penaliza la falta de suavidad sobre el grafo \cite{van2020survey,zhu2005semi}.

\begin{equation}
    \label{eq:minimizacionGrafos}
    \min_{f} \; \sum_{i \in L} \ell\bigl(f(\mathbf{x}_i), y_i\bigr) + \lambda \sum_{i,j} w_{ij}\, \ell_U\bigl(f(\mathbf{x}_i), f(\mathbf{x}_j)\bigr)
\end{equation}

En esta expresión, $L$ denota el conjunto de índices etiquetados y $U$ el de no etiquetados (con $L\cup U=\{1,\dots,n\}$). Además, $w_{ij}\ge 0$ cuantifica la similitud entre $\mathbf{x}_i$ y $\mathbf{x}_j$ (y, por convenio, $w_{ij}=0$ cuando no existe arista entre ambos nodos), mientras que $\ell_U$ es una penalización que favorece predicciones similares en nodos conectados.

\begin{rema}
    Nótese la similitud con el enfoque visto en los métodos de regularización de variedades. La diferencia aquí es que, en el planteamiento puramente transductivo, no se busca controlar la complejidad de un modelo definido sobre todo el espacio de entrada, sino obtener predicciones coherentes para las instancias del conjunto de entrenamiento; si se desea extrapolar a ejemplos nuevos, se pasa a formulaciones inductivas como las de la Sección \ref{sec: variedades}.
\end{rema}

Por su propia estructura, estos métodos tienen dos partes claramente diferenciadas: la \textbf{construcción del grafo} de similitud (definir vecindades y pesos) y la \textbf{inferencia sobre el grafo} (obtener la asignación de etiquetas o la función $f$ en los nodos) \cite{van2020survey}. Pasemos entonces a abordar cada una de estas partes.

\section{Construcción del grafo}

La construcción del grafo constituye la fase inicial de los métodos basados en grafos y, en la práctica, suele ser un componente determinante: la calidad del grafo de similitud condiciona de manera directa el rendimiento del método de inferencia posterior \cite{van2020survey}. Dado que los nodos representan instancias del conjunto de entrenamiento, construir el grafo equivale a decidir (i) qué pares de nodos se conectan mediante aristas y (ii) qué peso (similitud) se asigna a cada una de ellas.

\subsection{Construcción de la matriz de adyacencia}

El primer paso de la construcción de un grafo de similitud es definir una matriz de adyacencia $A$, donde cada elemento $a_{ij}$ indica la existencia de una arista entre los nodos $i$ y $j$. Existen varias formas de construir esta matriz, siendo las más comunes:

\begin{itemize}
    \item \textbf{Grafo completo}: se conecta cada par de nodos, es decir, $a_{ij}=1$ para todo $i\neq j$ \cite{zhu2009introduction}. Aunque conceptualmente sencillo, suele ser impracticable por su coste computacional y rara vez se utiliza en problemas de tamaño medio o grande.
    \item \textbf{Grafo de $\epsilon$-vecinos}: se conecta el nodo $i$ con el nodo $j$ si la distancia entre sus representaciones es menor que un umbral $\epsilon$, esto es, $a_{ij}=1$ si $d(\mathbf{x}_i,\mathbf{x}_j)<\epsilon$. Es un criterio simple, pero el grafo resultante depende fuertemente de la elección de $\epsilon$ y puede ser sensible a la escala de los datos \cite{van2020survey}.
    \item \textbf{Grafo de $k$ vecinos más cercanos (k-NN)}: se conecta el nodo $i$ con el nodo $j$ si $j$ pertenece al conjunto de los $k$ vecinos más cercanos de $i$. En general, este procedimiento produce un grafo dirigido (puede ocurrir que $j$ sea vecino de $i$ sin que $i$ lo sea de $j$). Para obtener un grafo no dirigido se suele \emph{simetrizar} la vecindad, por ejemplo conectando $i$ y $j$ si $i$ es vecino de $j$ \emph{o} $j$ es vecino de $i$, o bien exigiendo que la vecindad sea recíproca. Es uno de los enfoques más utilizados en la práctica, aunque también requiere seleccionar el parámetro $k$ \cite{van2020survey}.
    \item \textbf{Grafo mediante \textit{b-matching}}: mientras que los criterios anteriores son esencialmente locales, el \textit{b-matching} busca construir conexiones optimizando un objetivo global \cite{van2020survey}. De forma informal, se selecciona un conjunto de aristas que maximiza una medida de similitud agregada, sujeto a que cada nodo tenga a lo sumo $b$ conexiones. Puede producir grafos menos sensibles a parámetros locales, a costa de una implementación más compleja y un mayor coste computacional \cite{jebara2009graph}.
\end{itemize}

\subsection{Asignación de pesos}

Una vez definida la estructura del grafo, el siguiente paso es asignar pesos a las aristas. Estos pesos reflejan la similitud entre los nodos conectados, y su elección puede afectar significativamente el rendimiento del método de inferencia \cite{van2020survey}. Algunas formas comunes de asignar pesos son:

\begin{itemize}
    \item \textbf{Asignación gaussiana}: Una forma común de asignar pesos es usar una función gaussiana de la distancia entre los nodos, es decir, 
    \begin{equation}
        w_{ij} = \exp(-\frac{\|\mathbf{x}_i - \mathbf{x}_j\|^2}{2\sigma^2})
    \end{equation}
    donde $\sigma$ es un parámetro que controla la velocidad de decaimiento de la similitud \cite{zhu2009introduction}. Este método es popular porque produce pesos que decaen suavemente con la distancia, pero requiere elegir el parámetro $\sigma$ de manera adecuada.

    \item \textbf{Asignación gaussiana modificada}: Una variante de la anterior es usar una función gaussiana modificada que tenga en cuenta la densidad local de los datos, es decir,
    \begin{equation}
        w_{ij} = \exp(-\frac{\|\mathbf{x}_i - \mathbf{x}_j\|^2}{\max\{\sigma_i, \sigma_j\}^2})
    \end{equation}
    donde $\sigma_i$ y $\sigma_j$ son las distancias a los vecinos más cercanos de los nodos $i$ y $j$, respectivamente \cite{hein2006manifold}. 

    \item \textbf{Propagación lineal de vecinos}: Los métodos anteriores asignan $w_{ij}$ a partir de una función explícita de la distancia entre dos nodos. Sin embargo, existen alternativas que incorporan información de todo el vecindario. Un ejemplo es asignar pesos imponiendo que cada instancia pueda aproximarse como una combinación lineal de sus vecinos, minimizando un error de reconstrucción \cite{van2020survey}:
    \begin{equation}
        \min_{\{w_{ij}\}} \sum_{i} \left\| \mathbf{x}_i - \sum_{j \in N(i)} w_{ij} \, \mathbf{x}_j \right\|^2
    \end{equation}
    sujeto a la restricción $\sum_{j \in N(i)} w_{ij} = 1$ para cada nodo $i$, donde $N(i)$ denota el conjunto de vecinos de $i$ en el grafo. Este enfoque puede producir pesos mejor adaptados a la estructura local de los datos, aunque conlleva un mayor coste computacional \cite{van2020survey}.
\end{itemize}


\section{Inferencia sobre el grafo}

La construcción del grafo descrita en la sección anterior es necesaria siempre que se quiera explotar una noción de vecindad o geometría sobre el conjunto de entrenamiento; por lo que, también es un ingrediente central en los métodos de regularización de variedades (Sección \ref{sec: variedades}). La parte distintiva del enfoque transductivo es la \textbf{inferencia sobre el grafo}: dado el grafo de similitud y las etiquetas de un subconjunto de nodos, se trata de asignar etiquetas (o, más generalmente, una función de decisión) a los nodos no etiquetados, respetando la suavidad inducida por los pesos $w_{ij}$.

Una vez fijado el grafo, muchos métodos de inferencia pueden interpretarse como instancias del problema de optimización de la Ecuación \ref{eq:minimizacionGrafos}, que combina (i) un término de ajuste a las etiquetas conocidas y (ii) un término de regularización que penaliza variaciones bruscas de $f$ entre nodos similares. La elección de las pérdidas $\ell$ y $\ell_U$, del parámetro $\lambda$ y del algoritmo para resolver el problema determina el método concreto \cite{van2020survey}. A continuación se presentan dos enfoques clásicos.

\subsection{Mincut}

El método \textit{mincut} (o \textit{s-t mincut}) se plantea típicamente para clasificación binaria y busca un corte del grafo que separe los nodos etiquetados positivos de los negativos minimizando la suma de pesos de las aristas cortadas \cite{zhu2009introduction}. Una vez fijado el corte, los nodos no etiquetados heredan la etiqueta del lado del corte al que quedan conectados.

Atendiendo a la formulación de la Ecuación \ref{eq:minimizacionGrafos}, este método puede verse como la elección de una pérdida supervisada con peso infinito (equivalentemente, imponiendo como restricciones las etiquetas en $L$) junto con un término de suavidad cuadrático en el grafo. Una formulación habitual es

\begin{equation}
    \min_{f: f(\mathbf{x}_i)\in\{-1,1\}} \; \frac{1}{4} \sum_{i,j} w_{ij}\, \bigl(f(\mathbf{x}_i) - f(\mathbf{x}_j)\bigr)^2
    \quad \text{s.a.}\quad f(\mathbf{x}_i)=y_i,\; i\in L.
\end{equation}

El enfoque tiene limitaciones conocidas: en su forma básica está restringido a clasificación binaria y puede presentar cortes degenerados o no únicos, lo que se traduce en soluciones inestables o poco informativas para los nodos no etiquetados \cite{zhu2009introduction,van2020survey}.

\subsection{Funciones armónicas}

La idea de este método es relajar el problema de \textit{mincut} permitiendo que, en los nodos no etiquetados, la función $f$ tome valores continuos (por ejemplo, puntuaciones en $\mathbb{R}$), manteniendo las etiquetas conocidas fijadas en $L$. Gracias a esta relajación, el problema se reduce a la minimización de una función cuadrática, lo que conduce a una solución cerrada\footnote{Es decir, tiene una forma explícita que se calcula resolviendo un sistema de ecuaciones lineales} que, bajo condiciones suaves, es única y óptima \cite{zhu2009introduction}.

El nombre proviene de la conexión con las funciones armónicas: la solución satisface una condición de \emph{media ponderada} sobre el grafo. En particular, para cada nodo no etiquetado $i\in U$ se cumple \cite{van2020survey}:
\begin{equation}
    \sum_{j \in N(i)} w_{ij} (f(\mathbf{x}_i) - f(\mathbf{x}_j)) = 0
\end{equation}

equivalentemente, $f(\mathbf{x}_i)$ coincide con el promedio ponderado de sus vecinos, con pesos proporcionales a $w_{ij}$.


\backmatter

\printnoidxglossary[type=\acronymtype,title=Acrónimos]

\bibliographystyle{plain}
\bibliography{bibliografia}

\end{document}

